\documentclass[11pt,a4paper]{article}
\usepackage[utf8]{inputenc}
\usepackage[T1]{fontenc}
\usepackage[spanish,es-minimal,es-noshorthands]{babel}
\AtBeginDocument{\renewcommand{\sin}{\operatorname{sen}}\renewcommand{\arcsin}{\operatorname{arcsen}}}
\usepackage{lmodern}
\usepackage[sfdefault,lining]{FiraSans}
\usepackage[a4paper,top=2.3cm,bottom=2.2cm,left=1.8cm,right=1.8cm,headheight=1.3cm,headsep=0.4cm,footskip=1cm]{geometry}
\usepackage{amsmath,amssymb}
\usepackage{icomma}
\usepackage{graphicx}
\usepackage[table]{xcolor}
\usepackage{tikz}
\usepackage{fontawesome5}
\usepackage{ccicons}
\usepackage{fancyhdr}
\usepackage{lastpage}
\usepackage{tasks}
\usepackage{needspace}
\usepackage[most]{tcolorbox}
\usepackage{array}
\usepackage[hidelinks]{hyperref}
\AtBeginDocument{\hypersetup{pdftitle={Ficha de ejercicios de la unidad 6: Vectores y rectas. 4º ESO},pdfsubject={Matemáticas B 4º ESO},pdfkeywords={vectores, componentes, módulo, producto escalar, punto medio, distancia, puntos alineados, ecuaciones de la recta, pendiente, posición relativa, paralelas, perpendiculares, mediatriz}}}

\definecolor{navy}{HTML}{1B2748}
\definecolor{accent}{HTML}{E8603A}
\definecolor{teal}{HTML}{1F9E90}
\definecolor{soft}{HTML}{EEF1F8}
\definecolor{muted}{HTML}{5D6685}

\setlength{\parindent}{0pt}
\setlength{\parskip}{3pt}
\settasks{label=\alph*), label-format=\bfseries\color{navy}, label-width=1.3em, item-indent=1.8em, column-sep=0.6em, after-item-skip=5pt, before-skip=3pt, after-skip=4pt}

% ---------- capçalera i peu ----------
\newcommand{\llicencia}{\href{https://creativecommons.org/licenses/by-nc-sa/4.0/deed.es}{\ccbyncsa}\ \ Licencia Creative Commons Atribución-NoComercial-CompartirIgual 4.0 Internacional (CC BY-NC-SA 4.0)}
\pagestyle{fancy}
\fancyhf{}
\fancyhead[L]{\raisebox{-0.35\height}{\includegraphics[height=0.95cm]{logo-gva.png}}}
\fancyhead[C]{\small\color{muted}Ficha de ejercicios · Unidad 6: Vectores y rectas · 4\textsuperscript{o} ESO B}
\fancyhead[R]{\raisebox{-0.35\height}{\includegraphics[height=0.95cm]{logo-vermellar.png}}}
\fancyfoot[L]{\scriptsize\color{muted}\llicencia}
\fancyfoot[R]{\small\color{muted}\thepage}
\renewcommand{\headrulewidth}{0.4pt}
\renewcommand{\headrule}{\color{muted!40}\hrule height \headrulewidth}
\fancypagestyle{primera}{\fancyhf{}\renewcommand{\headrulewidth}{0pt}
  \fancyfoot[L]{\scriptsize\color{muted}\llicencia}
  \fancyfoot[R]{\small\color{muted}\thepage}}

% ---------- nivells i exercicis ----------
\newcommand{\estrelles}[1]{{\color{accent}\ifcase#1\or\faStar\faStar[regular]\faStar[regular]\or\faStar\faStar\faStar[regular]\or\faStar\faStar\faStar\fi}}
\newcounter{exer}
\newcommand{\ex}[1]{\par\needspace{4\baselineskip}\medskip\refstepcounter{exer}\noindent{\color{navy}\bfseries\theexer.}\ \ {\small\estrelles{#1}}\ \ }

\newcommand{\seccio}[2]{\par\needspace{16\baselineskip}\bigskip
  \begin{tcolorbox}[colback=navy,colframe=navy,coltext=white,arc=2mm,boxrule=0pt,left=3mm,right=3mm,top=1.2mm,bottom=1.2mm]
  {\large\bfseries #1\quad #2}
  \end{tcolorbox}\vspace{-2pt}}

\newcommand{\exf}[1]{\par\medskip\refstepcounter{exer}\noindent{\color{navy}\bfseries\theexer.}\ \ {\small\estrelles{#1}}\ \ }
\newcolumntype{C}[1]{>{\centering\arraybackslash}p{#1}}
% quadre de recordatori (només a la fitxa d'exercicis)
\newcommand{\recorda}[1]{\par\smallskip\begin{tcolorbox}[colback=accent!7,colframe=accent!55,arc=2mm,boxrule=0.5pt,left=2.5mm,right=2.5mm,top=1mm,bottom=1mm]\small{\color{accent}\faLightbulb[regular]}\ \textbf{Recuerda:}\ #1\end{tcolorbox}}
\usetikzlibrary{arrows.meta}
\usepackage{multicol}
\usepackage{enumitem}
\definecolor{violet}{HTML}{7C3AED}
\definecolor{orange}{HTML}{EA8A1F}
\newcommand{\res}[1]{#1}\newcommand{\rest}[1]{#1}

\def\fmneed{11}% que el títol d'una secció no quede sol al final del full
% =====================================================================
% Capçalera FentMates de les fitxes: un color per a cada curs.
% Els colors són els de la pàgina inicial (LEVELS a assets/js/materials.js);
% el logo ve de _eines/logo/pdf/ (python3 _eines/logo/genera.py).
%
% Ús, al final del preàmbul de la unitat:
%   % =====================================================================
% Capçalera FentMates de les fitxes: un color per a cada curs.
% Els colors són els de la pàgina inicial (LEVELS a assets/js/materials.js);
% el logo ve de _eines/logo/pdf/ (python3 _eines/logo/genera.py).
%
% Ús, al final del preàmbul de la unitat:
%   % =====================================================================
% Capçalera FentMates de les fitxes: un color per a cada curs.
% Els colors són els de la pàgina inicial (LEVELS a assets/js/materials.js);
% el logo ve de _eines/logo/pdf/ (python3 _eines/logo/genera.py).
%
% Ús, al final del preàmbul de la unitat:
%   \input{../comu/fentmates.tex}
%   \fmcurs{1eso}
% i, al començament del document:
%   \capcalera{Matemàtiques}{1r ESO}{Fitxa d'exercicis · Unitat 2}{Divisibilitat}
% Opcional abans de \input: \def\fmneed{14} (espai mínim abans de cada secció).
% =====================================================================
\makeatletter
\def\fm@logodir{../../logo/pdf}
\providecommand{\fmneed}{10}
\def\fm@cursnom{}

% clau del curs, color i color del text damunt del color
\newcommand{\fm@curs}[3]{\@namedef{fm@color@#1}{#2}\@namedef{fm@text@#1}{#3}}
\fm@curs{1eso}{E6B84A}{navy}
\fm@curs{2eso}{3B82F6}{white}
\fm@curs{3eso}{FF7A50}{white}
\fm@curs{4eso-a}{A78BFA}{white}
\fm@curs{4eso-b}{6D28D9}{white}
\fm@curs{1batx-mi}{9B1C1C}{white}
\fm@curs{1batx-ccss}{EF6B6B}{white}
\fm@curs{2batx-mii}{1E6B47}{white}
\fm@curs{2batx-ccss}{2BB3A3}{white}

\newcommand{\fmcurs}[1]{%
  \@ifundefined{fm@color@#1}{\PackageError{fentmates}{Curs desconegut: #1}{}}{}%
  \edef\fm@tmp{\noexpand\definecolor{curs}{HTML}{\@nameuse{fm@color@#1}}%
               \noexpand\colorlet{curstext}{\@nameuse{fm@text@#1}}}\fm@tmp
  \colorlet{cursfosc}{curs!72!black}%
  \def\fm@logo{\fm@logodir/vertical-#1.pdf}%
  \def\fm@marca{\fm@logodir/marca-#1.pdf}}

% ---------- capçalera de la primera pàgina ----------
% #1 matèria (petit), #2 curs, #3 tipus de document i unitat, #4 títol
\newcommand{\capcalera}[4]{%
  \gdef\fm@cursnom{#2}%
  \par\noindent
  \begin{tcolorbox}[colback=curs,colframe=curs,boxrule=0pt,arc=2.2mm,width=3.5cm,height=2.15cm,
      valign=center,left=3mm,right=2mm,top=1mm,bottom=1mm,nobeforeafter,box align=center]
    \color{curstext}{\footnotesize #1}\par\vspace{1pt}{\fontsize{19}{21}\selectfont\bfseries #2\par}%
  \end{tcolorbox}\hspace{5mm}%
  \begin{minipage}[c]{\dimexpr\linewidth-3.5cm-5mm-2.9cm\relax}
    {\color{cursfosc}\small\bfseries\MakeUppercase{#3}}\par\vspace{2pt}
    {\color{navy}\fontsize{20}{23}\selectfont\bfseries #4\par}
  \end{minipage}\hfill
  \begin{minipage}[c]{2.6cm}\raggedleft\includegraphics[width=2.6cm]{\fm@logo}\end{minipage}\par
  \vspace{2pt}}

% ---------- barres de secció del color del curs ----------
\renewcommand{\seccio}[2]{\par\needspace{\fmneed\baselineskip}\bigskip
  \begin{tcolorbox}[colback=curs,colframe=curs,coltext=curstext,arc=2mm,boxrule=0pt,left=3mm,right=3mm,top=1.2mm,bottom=1.2mm]
  {\large\bfseries #1\quad #2}
  \end{tcolorbox}\vspace{-2pt}}

% ---------- estrelles de dificultat del color del curs (a Batxillerat no se'n posen) ----------
\@ifundefined{estrelles}{}{%
  \renewcommand{\estrelles}[1]{{\color{curs}\ifcase#1\or\faStar\faStar[regular]\faStar[regular]\or\faStar\faStar\faStar[regular]\or\faStar\faStar\faStar\fi}}}

% ---------- capçalera i peu de les pàgines ----------
% primera pàgina: logos institucionals al peu
\fancypagestyle{primera}{\fancyhf{}\renewcommand{\headrulewidth}{0pt}
  \fancyfoot[L]{\raisebox{-0.32\height}{\includegraphics[height=0.78cm]{logo-gva.png}}\hspace{5mm}%
    \raisebox{-0.32\height}{\includegraphics[height=0.78cm]{logo-vermellar.png}}}
  \fancyfoot[R]{\scriptsize\color{muted}\href{https://creativecommons.org/licenses/by-nc-sa/4.0/deed.ca}{\ccbyncsa}\ \ CC BY-NC-SA 4.0 · fentmates.org\qquad\small\thepage}}
% la resta: el símbol a l'esquerra i el curs a la dreta, amb el títol al centre com abans
\fancyhead[L]{\raisebox{-0.3\height}{\includegraphics[height=0.62cm]{\fm@marca}}}
\fancyhead[R]{\small\bfseries\color{cursfosc}\fm@cursnom}
\renewcommand{\headrule}{\color{curs!45}\hrule height \headrulewidth}
\makeatother

%   \fmcurs{1eso}
% i, al començament del document:
%   \capcalera{Matemàtiques}{1r ESO}{Fitxa d'exercicis · Unitat 2}{Divisibilitat}
% Opcional abans de \input: \def\fmneed{14} (espai mínim abans de cada secció).
% =====================================================================
\makeatletter
\def\fm@logodir{../../logo/pdf}
\providecommand{\fmneed}{10}
\def\fm@cursnom{}

% clau del curs, color i color del text damunt del color
\newcommand{\fm@curs}[3]{\@namedef{fm@color@#1}{#2}\@namedef{fm@text@#1}{#3}}
\fm@curs{1eso}{E6B84A}{navy}
\fm@curs{2eso}{3B82F6}{white}
\fm@curs{3eso}{FF7A50}{white}
\fm@curs{4eso-a}{A78BFA}{white}
\fm@curs{4eso-b}{6D28D9}{white}
\fm@curs{1batx-mi}{9B1C1C}{white}
\fm@curs{1batx-ccss}{EF6B6B}{white}
\fm@curs{2batx-mii}{1E6B47}{white}
\fm@curs{2batx-ccss}{2BB3A3}{white}

\newcommand{\fmcurs}[1]{%
  \@ifundefined{fm@color@#1}{\PackageError{fentmates}{Curs desconegut: #1}{}}{}%
  \edef\fm@tmp{\noexpand\definecolor{curs}{HTML}{\@nameuse{fm@color@#1}}%
               \noexpand\colorlet{curstext}{\@nameuse{fm@text@#1}}}\fm@tmp
  \colorlet{cursfosc}{curs!72!black}%
  \def\fm@logo{\fm@logodir/vertical-#1.pdf}%
  \def\fm@marca{\fm@logodir/marca-#1.pdf}}

% ---------- capçalera de la primera pàgina ----------
% #1 matèria (petit), #2 curs, #3 tipus de document i unitat, #4 títol
\newcommand{\capcalera}[4]{%
  \gdef\fm@cursnom{#2}%
  \par\noindent
  \begin{tcolorbox}[colback=curs,colframe=curs,boxrule=0pt,arc=2.2mm,width=3.5cm,height=2.15cm,
      valign=center,left=3mm,right=2mm,top=1mm,bottom=1mm,nobeforeafter,box align=center]
    \color{curstext}{\footnotesize #1}\par\vspace{1pt}{\fontsize{19}{21}\selectfont\bfseries #2\par}%
  \end{tcolorbox}\hspace{5mm}%
  \begin{minipage}[c]{\dimexpr\linewidth-3.5cm-5mm-2.9cm\relax}
    {\color{cursfosc}\small\bfseries\MakeUppercase{#3}}\par\vspace{2pt}
    {\color{navy}\fontsize{20}{23}\selectfont\bfseries #4\par}
  \end{minipage}\hfill
  \begin{minipage}[c]{2.6cm}\raggedleft\includegraphics[width=2.6cm]{\fm@logo}\end{minipage}\par
  \vspace{2pt}}

% ---------- barres de secció del color del curs ----------
\renewcommand{\seccio}[2]{\par\needspace{\fmneed\baselineskip}\bigskip
  \begin{tcolorbox}[colback=curs,colframe=curs,coltext=curstext,arc=2mm,boxrule=0pt,left=3mm,right=3mm,top=1.2mm,bottom=1.2mm]
  {\large\bfseries #1\quad #2}
  \end{tcolorbox}\vspace{-2pt}}

% ---------- estrelles de dificultat del color del curs (a Batxillerat no se'n posen) ----------
\@ifundefined{estrelles}{}{%
  \renewcommand{\estrelles}[1]{{\color{curs}\ifcase#1\or\faStar\faStar[regular]\faStar[regular]\or\faStar\faStar\faStar[regular]\or\faStar\faStar\faStar\fi}}}

% ---------- capçalera i peu de les pàgines ----------
% primera pàgina: logos institucionals al peu
\fancypagestyle{primera}{\fancyhf{}\renewcommand{\headrulewidth}{0pt}
  \fancyfoot[L]{\raisebox{-0.32\height}{\includegraphics[height=0.78cm]{logo-gva.png}}\hspace{5mm}%
    \raisebox{-0.32\height}{\includegraphics[height=0.78cm]{logo-vermellar.png}}}
  \fancyfoot[R]{\scriptsize\color{muted}\href{https://creativecommons.org/licenses/by-nc-sa/4.0/deed.ca}{\ccbyncsa}\ \ CC BY-NC-SA 4.0 · fentmates.org\qquad\small\thepage}}
% la resta: el símbol a l'esquerra i el curs a la dreta, amb el títol al centre com abans
\fancyhead[L]{\raisebox{-0.3\height}{\includegraphics[height=0.62cm]{\fm@marca}}}
\fancyhead[R]{\small\bfseries\color{cursfosc}\fm@cursnom}
\renewcommand{\headrule}{\color{curs!45}\hrule height \headrulewidth}
\makeatother

%   \fmcurs{1eso}
% i, al començament del document:
%   \capcalera{Matemàtiques}{1r ESO}{Fitxa d'exercicis · Unitat 2}{Divisibilitat}
% Opcional abans de \input: \def\fmneed{14} (espai mínim abans de cada secció).
% =====================================================================
\makeatletter
\def\fm@logodir{../../logo/pdf}
\providecommand{\fmneed}{10}
\def\fm@cursnom{}

% clau del curs, color i color del text damunt del color
\newcommand{\fm@curs}[3]{\@namedef{fm@color@#1}{#2}\@namedef{fm@text@#1}{#3}}
\fm@curs{1eso}{E6B84A}{navy}
\fm@curs{2eso}{3B82F6}{white}
\fm@curs{3eso}{FF7A50}{white}
\fm@curs{4eso-a}{A78BFA}{white}
\fm@curs{4eso-b}{6D28D9}{white}
\fm@curs{1batx-mi}{9B1C1C}{white}
\fm@curs{1batx-ccss}{EF6B6B}{white}
\fm@curs{2batx-mii}{1E6B47}{white}
\fm@curs{2batx-ccss}{2BB3A3}{white}

\newcommand{\fmcurs}[1]{%
  \@ifundefined{fm@color@#1}{\PackageError{fentmates}{Curs desconegut: #1}{}}{}%
  \edef\fm@tmp{\noexpand\definecolor{curs}{HTML}{\@nameuse{fm@color@#1}}%
               \noexpand\colorlet{curstext}{\@nameuse{fm@text@#1}}}\fm@tmp
  \colorlet{cursfosc}{curs!72!black}%
  \def\fm@logo{\fm@logodir/vertical-#1.pdf}%
  \def\fm@marca{\fm@logodir/marca-#1.pdf}}

% ---------- capçalera de la primera pàgina ----------
% #1 matèria (petit), #2 curs, #3 tipus de document i unitat, #4 títol
\newcommand{\capcalera}[4]{%
  \gdef\fm@cursnom{#2}%
  \par\noindent
  \begin{tcolorbox}[colback=curs,colframe=curs,boxrule=0pt,arc=2.2mm,width=3.5cm,height=2.15cm,
      valign=center,left=3mm,right=2mm,top=1mm,bottom=1mm,nobeforeafter,box align=center]
    \color{curstext}{\footnotesize #1}\par\vspace{1pt}{\fontsize{19}{21}\selectfont\bfseries #2\par}%
  \end{tcolorbox}\hspace{5mm}%
  \begin{minipage}[c]{\dimexpr\linewidth-3.5cm-5mm-2.9cm\relax}
    {\color{cursfosc}\small\bfseries\MakeUppercase{#3}}\par\vspace{2pt}
    {\color{navy}\fontsize{20}{23}\selectfont\bfseries #4\par}
  \end{minipage}\hfill
  \begin{minipage}[c]{2.6cm}\raggedleft\includegraphics[width=2.6cm]{\fm@logo}\end{minipage}\par
  \vspace{2pt}}

% ---------- barres de secció del color del curs ----------
\renewcommand{\seccio}[2]{\par\needspace{\fmneed\baselineskip}\bigskip
  \begin{tcolorbox}[colback=curs,colframe=curs,coltext=curstext,arc=2mm,boxrule=0pt,left=3mm,right=3mm,top=1.2mm,bottom=1.2mm]
  {\large\bfseries #1\quad #2}
  \end{tcolorbox}\vspace{-2pt}}

% ---------- estrelles de dificultat del color del curs (a Batxillerat no se'n posen) ----------
\@ifundefined{estrelles}{}{%
  \renewcommand{\estrelles}[1]{{\color{curs}\ifcase#1\or\faStar\faStar[regular]\faStar[regular]\or\faStar\faStar\faStar[regular]\or\faStar\faStar\faStar\fi}}}

% ---------- capçalera i peu de les pàgines ----------
% primera pàgina: logos institucionals al peu
\fancypagestyle{primera}{\fancyhf{}\renewcommand{\headrulewidth}{0pt}
  \fancyfoot[L]{\raisebox{-0.32\height}{\includegraphics[height=0.78cm]{logo-gva.png}}\hspace{5mm}%
    \raisebox{-0.32\height}{\includegraphics[height=0.78cm]{logo-vermellar.png}}}
  \fancyfoot[R]{\scriptsize\color{muted}\href{https://creativecommons.org/licenses/by-nc-sa/4.0/deed.ca}{\ccbyncsa}\ \ CC BY-NC-SA 4.0 · fentmates.org\qquad\small\thepage}}
% la resta: el símbol a l'esquerra i el curs a la dreta, amb el títol al centre com abans
\fancyhead[L]{\raisebox{-0.3\height}{\includegraphics[height=0.62cm]{\fm@marca}}}
\fancyhead[R]{\small\bfseries\color{cursfosc}\fm@cursnom}
\renewcommand{\headrule}{\color{curs!45}\hrule height \headrulewidth}
\makeatother

\fmcurs{4eso-b}

\begin{document}
\thispagestyle{primera}

\capcalera{Matemáticas B}{4\textsuperscript{o} ESO}{Ficha de ejercicios · Unidad 6}{\raggedright Vectores y rectas}
\vspace{10pt}
Nombre y apellidos: \rule{6.6cm}{0.4pt}\hfill Grupo: \rule{1.6cm}{0.4pt}\hfill Fecha: \rule{2.2cm}{0.4pt}
\vspace{8pt}

\begin{tcolorbox}[colback=soft,colframe=soft,arc=2mm,left=3mm,right=3mm,top=1.5mm,bottom=1.5mm]
\small Ejercicios graduados por niveles de dificultad:\quad \estrelles{1}\ básico \qquad \estrelles{2}\ medio \qquad \estrelles{3}\ avanzado.\\[2pt] Resuélvelos en el cuaderno mostrando todos los pasos. Haz un dibujo aproximado en unos ejes de coordenadas siempre que te ayude: te permite detectar errores. Da los resultados de manera exacta (fracciones y raíces) y, si hace falta, también aproximados. La ficha recoge todos los ejercicios de las fichas del departamento y añade otros nuevos.
\end{tcolorbox}

% =====================================================
\seccio{1}{Vectores}

\exf{1} Calcula las componentes del vector $\overrightarrow{AB}$ conociendo su origen $A$ y su extremo $B$:\begin{tasks}(2)
  \task $A(2,\,5)$, \ $B(4,\,-1)$
  \task $A(-3,\,2)$, \ $B(1,\,6)$
  \task $A(0,\,-4)$, \ $B(-5,\,-1)$
  \task $A(3,\,3)$, \ $B(3,\,-2)$
\end{tasks}

\ex{1} Calcula el módulo de los vectores siguientes:\begin{tasks}(3)
  \task $\vec a=(6,\,8)$
  \task $\vec b=(-5,\,12)$
  \task $\vec c=(2,\,-3)$
  \task $\vec d=(-4,\,-4)$
  \task $\vec e=(0,\,-7)$
  \task $\vec f=(8,\,-6)$
\end{tasks}

\ex{1} Sean los puntos $A(3,-2)$, $B(-4,1)$ y $C(5,2)$. Halla:\begin{enumerate}[label=\alph*),leftmargin=*,itemsep=2pt,topsep=2pt]\item Los vectores $\overrightarrow{AB}$ y $\overrightarrow{BC}$.\item El módulo del vector $\overrightarrow{AB}$.\item La distancia del punto $A$ al punto $C$.\end{enumerate}

\ex{1} Calcula el punto que falta en cada caso:\begin{enumerate}[label=\alph*),leftmargin=*,itemsep=2pt,topsep=2pt]\item El extremo $B$, si $\overrightarrow{AB}=\vec v$ tiene el origen en $A$: \ $\vec v=(4,\,-1)$, $A(-3,\,2)$; \ $\vec v=(-2,\,-5)$, $A(1,\,3)$; \ $\vec v=(6,\,0)$, $A(-4,\,-4)$; \ $\vec v=(-3,\,7)$, $A(5,\,-6)$.\item El origen $A$, si $\overrightarrow{AB}=\vec v$ tiene el extremo en $B$: \ $\vec v=(-2,\,4)$, $B(1,\,1)$; \ $\vec v=(5,\,-3)$, $B(0,\,0)$; \ $\vec v=(-1,\,-6)$, $B(-4,\,2)$; \ $\vec v=(3,\,3)$, $B(2,\,-5)$.\end{enumerate}

\ex{2} El vector $\vec v=(4,-3)$ se dibuja a partir de diferentes puntos. Escribe tres pares de puntos origen-extremo que lo representen: uno con el origen en $O(0,0)$, otro con el origen en $(-1,5)$ y otro con el extremo en $O(0,0)$.

\ex{1} Observa los vectores del dibujo.\begin{enumerate}[label=\alph*),leftmargin=*,itemsep=2pt,topsep=2pt]\item Indica el origen, el extremo y las componentes de cada vector, y calcula su módulo.\item ¿Qué vectores son equipolentes? ¿Cuál es el opuesto de $\vec a$?\end{enumerate}

\begin{center}\begin{tikzpicture}[scale=0.45,baseline=(current bounding box.center)]
\draw[muted!25,very thin] (-6,-5) grid (6,6);
\draw[-latex,muted] (-6.3,0)--(6.5,0) node[right,font=\scriptsize]{$x$}; \draw[-latex,muted] (0,-5.3)--(0,6.5) node[above,font=\scriptsize]{$y$};
\foreach \i in {-6,-4,...,6}{\ifnum\i=0\else\node[font=\tiny,text=muted,below] at (\i,0) {$\i$};\fi}
\foreach \i in {-4,-2,...,6}{\ifnum\i=0\else\node[font=\tiny,text=muted,left] at (0,\i) {$\i$};\fi}
\begin{scope}\clip (-6,-5) rectangle (6,6);
\end{scope}
\draw[-{Latex[length=2.2mm]},accent,line width=1.1pt] (-5,1)--(-2,3);
\node[accent,font=\small,above left] at (-3.5,2) {$\vec a$};
\draw[-{Latex[length=2.2mm]},teal,line width=1.1pt] (1,3)--(4,5);
\node[teal,font=\small,above left] at (2.5,4) {$\vec b$};
\draw[-{Latex[length=2.2mm]},navy,line width=1.1pt] (-4,-1)--(-4,-4);
\node[navy,font=\small,above left] at (-4,-2.5) {$\vec c$};
\draw[-{Latex[length=2.2mm]},violet,line width=1.1pt] (5,-1)--(2,-3);
\node[violet,font=\small,above left] at (3.5,-2) {$\vec d$};
\draw[-{Latex[length=2.2mm]},orange,line width=1.1pt] (-1,-3)--(3,-4);
\node[orange,font=\small,above left] at (1,-3.5) {$\vec e$};
\end{tikzpicture}\end{center}


\ex{2} Representa el triángulo de vértices $A(-2,1)$, $B(3,4)$ y $C(2,-2)$ y calcula las componentes de los vectores $\overrightarrow{AB}$, $\overrightarrow{BC}$ y $\overrightarrow{CA}$. ¿Cuánto vale $\overrightarrow{AB}+\overrightarrow{BC}+\overrightarrow{CA}$? ¿Por qué?


% =====================================================
\seccio{2}{Operaciones con vectores}

\exf{1} Dados los vectores $\vec u=(2,-3)$ y $\vec v=(-1,4)$, calcula:\begin{tasks}(3)
  \task $\vec u+\vec v$
  \task $\vec u-\vec v$
  \task $3\vec u$
  \task $-2\vec v$
  \task $2\vec u+3\vec v$
  \task $\tfrac12\vec u-\vec v$
\end{tasks}

\ex{1} Con $\vec a=(1,-4)$, $\vec b=(-2,3)$ y $\vec c=(5,0)$, opera:\begin{tasks}(2)
  \task $\vec a+\vec b+\vec c$
  \task $2\vec a-\vec b+\vec c$
  \task $-\vec a+3\vec b-2\vec c$
  \task $\tfrac15\vec c-\tfrac12\vec b$
\end{tasks}

\ex{2} Dibuja los vectores $\vec u=(3,1)$ y $\vec v=(-1,2)$ con origen en $O(0,0)$ y, gráficamente, $\vec u+\vec v$ (regla del paralelogramo), $\vec u-\vec v$ y $2\vec u$. Comprueba analíticamente los resultados.

\ex{1} Averigua si son paralelos los vectores $\vec u$ y $\vec v$:\begin{tasks}(3)
  \task $\vec u=(2,\,-6)$, $\vec v=(-1,\,3)$
  \task $\vec u=(4,\,6)$, $\vec v=(6,\,9)$
  \task $\vec u=(3,\,5)$, $\vec v=(5,\,3)$
  \task $\vec u=(-2,\,8)$, $\vec v=(1,\,4)$
  \task $\vec u=(0,\,5)$, $\vec v=(0,\,-2)$
  \task $\vec u=(10,\,-4)$, $\vec v=(-15,\,6)$
\end{tasks}

\ex{2} Calcula $x$ en cada caso para que los vectores sean paralelos:\begin{tasks}(3)
  \task $\vec u=(x,3)$, $\vec v=(12,x)$
  \task $\vec a=(2,-5)$, $\vec b=(x,10)$
  \task $\vec c=(x,x+1)$, $\vec d=(3,4)$
\end{tasks}

\ex{2} Escribe un vector $\vec u$ con la misma dirección que $\vec v=(-3,9)$, sentido contrario y módulo la tercera parte del de $\vec v$. ¿Qué operación hay que aplicar a $\vec u$ para volver a obtener $\vec v$?

\ex{3} Halla el vector $\vec x$ que cumple $3\vec x+\vec u=2\vec v$, sabiendo que $\vec u=(1,-5)$ y $\vec v=(2,-1)$.


% =====================================================
\seccio{3}{Producto escalar. Vectores perpendiculares}

\exf{1} Calcula el producto escalar $\vec u\cdot\vec v$ e indica si los vectores forman un ángulo agudo, recto u obtuso:\begin{tasks}(2)
  \task $\vec u=(3,\,2)$, $\vec v=(1,\,-4)$
  \task $\vec u=(-2,\,5)$, $\vec v=(4,\,-1)$
  \task $\vec u=(6,\,-3)$, $\vec v=(1,\,2)$
  \task $\vec u=(0,\,3)$, $\vec v=(-5,\,7)$
\end{tasks}

\ex{1} Haz las siguientes operaciones con vectores sabiendo que $A(8,-4)$, $B(1,5)$, $C(7,-2)$ y $D(3,-8)$:\begin{tasks}(3)
  \task $\overrightarrow{AB}+\overrightarrow{CD}$
  \task $\overrightarrow{AB}-3\overrightarrow{CD}$
  \task $2\overrightarrow{AB}\cdot\overrightarrow{CD}$
  \task $5\overrightarrow{AB}\cdot(-\overrightarrow{DA})$
  \task $|\overrightarrow{AB}|$
  \task $|\overrightarrow{CD}|$
\end{tasks}

\ex{1} Determina si estos pares de vectores son paralelos, perpendiculares o ninguna de las dos cosas:\begin{tasks}(3)
  \task $\vec u=(2,\,1)$, $\vec v=(-1,\,2)$
  \task $\vec u=(-6,\,4)$, $\vec v=(3,\,-2)$
  \task $\vec u=(5,\,2)$, $\vec v=(2,\,5)$
  \task $\vec u=(4,\,-6)$, $\vec v=(9,\,6)$
  \task $\vec u=(1,\,1)$, $\vec v=(-3,\,-3)$
  \task $\vec u=(7,\,0)$, $\vec v=(0,\,-2)$
\end{tasks}

\ex{2} Escribe dos vectores perpendiculares a $\vec v=(5,-2)$ y dos paralelos, uno con el mismo sentido y otro con sentido contrario. Demuestra que, para cualquier vector $\vec v=(v_1,v_2)$, el vector $\vec w=(-v_2,v_1)$ es perpendicular a $\vec v$.

\ex{2} Calcula $x$ para que los vectores sean perpendiculares:\begin{tasks}(3)
  \task $\vec u=(3,x)$, $\vec v=(x+1,-2)$
  \task $\vec a=(x,-4)$, $\vec b=(6,3)$
  \task $\vec c=(x,2)$, $\vec d=(x,-8)$
\end{tasks}

\ex{2} Representa los vectores $\vec a=(2,-1)$, $\vec b=(1,2)$ y $\vec c=(-4,2)$. ¿Qué relación hay entre cada par?

\ex{3} Halla los vectores perpendiculares a $\vec v=(3,4)$ que tienen módulo 10.


% =====================================================
\seccio{4}{Aplicaciones de los vectores}

\exf{1} Calcula la distancia entre los puntos $A(4,0)$ y $B(-1,1)$. Representa el vector $\overrightarrow{AB}$. Escribe un vector perpendicular y uno paralelo a $\overrightarrow{AB}$.

\ex{1} Calcula la distancia entre los puntos $A$ y $B$:\begin{tasks}(3)
  \task $A(1,\,2)$, $B(4,\,6)$
  \task $A(-3,\,5)$, $B(2,\,-7)$
  \task $A(0,\,-2)$, $B(-4,\,-2)$
  \task $A(-1,\,-1)$, $B(2,\,2)$
  \task $A(3,\,-4)$, $B(-2,\,1)$
  \task $A(-6,\,0)$, $B(0,\,8)$
\end{tasks}

\ex{1} Calcula el punto medio del segmento de extremos $P$ y $Q$ (el último es de la ficha del departamento):\begin{tasks}(3)
  \task $P(2,\,6)$, $Q(8,\,0)$
  \task $P(-3,\,4)$, $Q(5,\,-2)$
  \task $P(1,\,-5)$, $Q(-2,\,2)$
  \task $P(0,\,7)$, $Q(-6,\,-3)$
  \task $P(2,\,1)$, $Q(-4,\,3)$
\end{tasks}

\ex{1} $M$ es el punto medio del segmento $AB$. Calcula el extremo $B$:\begin{tasks}(2)
  \task $A(1,\,4)$, $M(3,\,1)$
  \task $A(-2,\,0)$, $M(0,\,3)$
  \task $A(4,\,-3)$, $M(1,\,-1)$
  \task $A(-5,\,2)$, $M(-1,\,-2)$
\end{tasks}

\ex{1} Calcula el punto simétrico $A'$ del punto $A(-1,0)$ respecto de $B(2,-8)$. Calcula también el simétrico de $P(3,-2)$ respecto de $Q(1,1)$.

\ex{2} Puntos alineados:\begin{enumerate}[label=\alph*),leftmargin=*,itemsep=2pt,topsep=2pt]\item Determina si los puntos $A(3,1)$, $B(5,2)$ y $C(1,0)$ están alineados.\item Calcula el valor de $k$ para que los puntos $A(1,1)$, $B(0,3)$ y $C(2,k)$ estén alineados.\item Comprueba que los puntos $P(-1,2)$, $Q(1,5)$ y $R(4,7)$ no están alineados y que, por tanto, forman un triángulo.\end{enumerate}

\ex{2} Paralelogramos:\begin{enumerate}[label=\alph*),leftmargin=*,itemsep=2pt,topsep=2pt]\item Calcula el cuarto vértice $B$ del paralelogramo $ABCD$ si $A(2,1)$, $C(10,-1)$ y $D(7,-3)$.\item Calcula el punto $D$ para que $A(-2,1)$, $B(1,3)$, $C(4,0)$ y $D$ sean los vértices (en este orden) de un paralelogramo.\end{enumerate}

\ex{2} Comprueba mediante vectores que el triángulo de vértices $A(-4,2)$, $B(5,-1)$ y $C(-2,8)$ es rectángulo.

\ex{2} Los puntos $A(4,-2)$, $B(5,2)$ y $C(9,1)$ son vértices de un cuadrado. Halla el otro vértice.

\ex{2} Verifica que el cuadrilátero de vértices $A(1,1)$, $B(4,2)$, $C(3,5)$ y $D(0,4)$ es un cuadrado y calcula su perímetro y su área.

\ex{2} Los vértices de un rombo son $A(2,5)$, $B(0,2)$, $C(2,-1)$ y $D(4,2)$. Comprueba que todos los lados miden lo mismo y calcula la longitud de las diagonales y el área del rombo.

\ex{2} Calcula los vectores que forman los lados del triángulo de vértices $A(-2,-1)$, $B(4,2)$ y $C(1,5)$, y su perímetro. ¿Qué tipo de triángulo es según sus lados?

\ex{2} Circunferencias:\begin{enumerate}[label=\alph*),leftmargin=*,itemsep=2pt,topsep=2pt]\item Los puntos $P(-3,4)$ y $Q(5,-2)$ son los extremos de un diámetro de una circunferencia. Calcula su centro y su radio.\item Dos circunferencias del mismo radio son tangentes exteriores y tienen los centros en $A(-2,3)$ y $B(6,-1)$. Halla el punto de tangencia.\end{enumerate}

\ex{3} Halla el punto $P$ del eje de abscisas que está a la misma distancia de $A(1,3)$ que de $B(5,1)$.

\ex{3} En el triángulo de vértices $A(0,0)$, $B(6,2)$ y $C(2,6)$, sean $M$ y $N$ los puntos medios de los lados $AC$ y $BC$. Comprueba que el segmento $MN$ es paralelo al lado $AB$ y que mide la mitad.


% =====================================================
\seccio{5}{Ecuaciones vectorial, paramétricas y continua}

\exf{1} Escribe las ecuaciones vectorial, paramétricas y continua de la recta que pasa por el punto $A$ y tiene vector director $\vec v$:\begin{tasks}(2)
  \task $A(2,\,-1)$, $\vec v=(3,\,4)$
  \task $A(0,\,5)$, $\vec v=(-2,\,1)$
  \task $A(-4,\,-3)$, $\vec v=(1,\,-5)$
  \task $A(6,\,0)$, $\vec v=(-3,\,-3)$
\end{tasks}

\ex{1} Escribe las ecuaciones vectorial, paramétricas y continua de la recta que pasa por los puntos $A$ y $B$:\begin{tasks}(2)
  \task $A(1,\,2)$, $B(4,\,-1)$
  \task $A(-2,\,0)$, $B(3,\,4)$
  \task $A(5,\,-3)$, $B(-1,\,-6)$
  \task $A(0,\,0)$, $B(-4,\,6)$
\end{tasks}

\ex{1} Escribe la ecuación vectorial y las paramétricas de las rectas representadas. ¿Cuál de ellas no se puede escribir en forma continua? ¿Por qué?

\begin{center}\begin{tikzpicture}[scale=0.42,baseline=(current bounding box.center)]
\draw[muted!25,very thin] (-6,-6) grid (6,6);
\draw[-latex,muted] (-6.3,0)--(6.5,0) node[right,font=\scriptsize]{$x$}; \draw[-latex,muted] (0,-6.3)--(0,6.5) node[above,font=\scriptsize]{$y$};
\foreach \i in {-6,-4,...,6}{\ifnum\i=0\else\node[font=\tiny,text=muted,below] at (\i,0) {$\i$};\fi}
\foreach \i in {-6,-4,...,6}{\ifnum\i=0\else\node[font=\tiny,text=muted,left] at (0,\i) {$\i$};\fi}
\begin{scope}\clip (-6,-6) rectangle (6,6);
\draw[accent,thick] (-6,-4.33333)--(6,3.66667);
\draw[teal,thick] (-6,16)--(6,-8);
\draw[navy,thick] (-6,5.33333)--(6,1.33333);
\draw[violet,thick] (4,-6)--(4,6);
\end{scope}
\node[accent,font=\small\itshape,fill=white,inner sep=1pt] at (-5.3,-3.3) {$r$};
\node[teal,font=\small\itshape,fill=white,inner sep=1pt] at (3.5,-2.1) {$s$};
\node[navy,font=\small\itshape,fill=white,inner sep=1pt] at (-5.4,4.3) {$t$};
\node[violet,font=\small\itshape,fill=white,inner sep=1pt] at (4.4,-5.3) {$u$};
\end{tikzpicture}\end{center}


\ex{2} Comprueba si el punto $P(3,-1)$ pertenece a cada una de estas rectas:\begin{tasks}[after-item-skip=6pt](2)
  \task $(x,y)=(1,3)+\lambda(1,-2)$
  \task $\left\{\begin{array}{l}x=5-2\lambda\\ y=2+\lambda\end{array}\right.$
  \task $\dfrac{x-7}{4}=\dfrac{y+3}{-2}$
  \task $\dfrac{x+1}{2}=\dfrac{y-1}{3}$
\end{tasks}

\ex{2} Dada la recta $r:\left\{\begin{array}{l}x=3+2\lambda\\ y=1-4\lambda\end{array}\right.$\begin{enumerate}[label=\alph*),leftmargin=*,itemsep=2pt,topsep=2pt]\item Halla tres puntos de $r$.\item Calcula $k$ para que el punto $(k,-7)$ esté en $r$.\item Escribe la ecuación continua de $r$ y comprueba con ella el resultado del apartado anterior.\end{enumerate}

\ex{2} Escribe las ecuaciones de las rectas paralelas a los ejes que pasan por $P(-2,5)$: en forma vectorial y de manera sencilla. ¿Cuáles son las ecuaciones de los ejes de coordenadas?


% =====================================================
\seccio{6}{Pendiente. Ecuaciones punto-pendiente, explícita y general}

\exf{1} Calcula la pendiente de la recta que pasa por los puntos $A$ y $B$. ¿Es creciente o decreciente?\begin{tasks}(2)
  \task $A(1,\,2)$, $B(3,\,8)$
  \task $A(-2,\,4)$, $B(2,\,2)$
  \task $A(0,\,-3)$, $B(5,\,-3)$
  \task $A(4,\,1)$, $B(4,\,6)$
\end{tasks}

\ex{1} Escribe la ecuación punto-pendiente y la explícita de la recta que tiene pendiente $m$ y pasa por el punto $A$:\begin{tasks}(2)
  \task $m=4$, $A(1,\,-2)$
  \task $m=-3$, $A(0,\,5)$
  \task $m=\frac{1}{2}$, $A(-4,\,1)$
  \task $m=-\frac{2}{3}$, $A(3,\,0)$
\end{tasks}

\ex{1} La ecuación de la recta $r$ es $y=-x+2$.\begin{enumerate}[label=\alph*),leftmargin=*,itemsep=2pt,topsep=2pt]\item ¿Cuál es el valor de la pendiente?\item Determina las coordenadas de uno de sus vectores directores.\item Halla dos puntos de la recta y dibújala.\item ¿El punto $A(-1,4)$ pertenece a esta recta?\end{enumerate}

\ex{1} La recta $r$ pasa por el punto $P(2,2)$ y tiene vector director $\vec v=(-4,2)$. Halla todas sus ecuaciones.

\ex{1} Halla las ecuaciones de la recta que pasa por los puntos $P(-3,-4)$ y $Q(-1,-2)$.

\ex{1} La recta $r$ pasa por el punto $P(5,-1)$ y tiene pendiente 2. Halla todas las ecuaciones de la recta.

\ex{1} Escribe la ecuación vectorial de la recta $r$ que pasa por los puntos $A(1,0)$ y $B(3,6)$, y deduce todas las demás ecuaciones.\label{ex:r7}

\ex{2} Escribe la ecuación de las rectas siguientes en todas las formas posibles:\begin{tasks}[after-item-skip=6pt](3)
  \task $\left\{\begin{array}{l}x=2-t\\ y=-3+2t\end{array}\right.$
  \task $(x,y)=(0,3)+t(2,1)$
  \task $y=3x-1$
  \task $y-3=3(x-5)$
  \task $2x+y-5=0$
\end{tasks}

\ex{2} Deduce si el punto $M(-2,1)$ pertenece o no a cada una de las rectas siguientes:\begin{tasks}[after-item-skip=6pt](3)
  \task $(x,y)=(3,2)+t(4,-3)$
  \task $\dfrac{x-5}{2}=\dfrac{3-y}{4}$
  \task $3x+5y+30=0$
  \task $y=\dfrac{x}{2}+2$
  \task $y+2=3(x+5)$
\end{tasks}

\ex{1} Indica un punto y un vector de las rectas:\begin{tasks}[after-item-skip=6pt](2)
  \task $(x,y)=(2,4)+\lambda(5,-3)$
  \task $\dfrac{x-3}{3}=\dfrac{y+2}{2}$
  \task $\left\{\begin{array}{l}x=-1+9\lambda\\ y=-8-6\lambda\end{array}\right.$
  \task $2x-3y+8=0$
\end{tasks}

\ex{1} Calcula la pendiente, la ordenada en el origen, un vector director y un vector normal de las rectas:\begin{tasks}(2)
  \task $3x-4y+12=0$
  \task $2x+y-5=0$
  \task $x-3=0$
  \task $5y+10=0$
\end{tasks}

\ex{2} Calcula la ecuación general de la recta:\begin{tasks}(2)
  \task Que pasa por $A(2,-3)$ y $B(5,1)$.
  \task Que pasa por $A(-1,4)$ y $B(3,4)$.
  \task Que pasa por $A(0,-2)$ con $\vec v=(3,-1)$.
  \task Que pasa por $A(4,1)$ con $\vec n=(2,-5)$.
\end{tasks}

\ex{1} Calcula la ecuación explícita de la recta que pasa por el punto $A$ y tiene vector director $\vec v$:\begin{tasks}(2)
  \task $A(2,\,-1)$, $\vec v=(4,\,2)$
  \task $A(-3,\,5)$, $\vec v=(-2,\,6)$
  \task $A(0,\,4)$, $\vec v=(3,\,-2)$
  \task $A(-1,\,-2)$, $\vec v=(5,\,0)$
\end{tasks}

\ex{2} Calcula la ecuación explícita de la recta en cada caso:\begin{enumerate}[label=\alph*),leftmargin=*,itemsep=2pt,topsep=2pt]\item La pendiente es $-\frac32$ y corta el eje $OX$ en el punto de abscisa $x=4$.\item Es paralela al eje $OX$ y su ordenada en el origen es $-2$.\item Pasa por el origen de coordenadas y por el punto $(2,-5)$.\end{enumerate}

\ex{3} La recta $r:3x+4y-12=0$ corta los ejes de coordenadas en los puntos $P$ y $Q$. Calcúlalos y halla el área y el perímetro del triángulo que forman con el origen.


% =====================================================
\seccio{7}{Posición relativa de dos rectas}

\exf{1} Calcula el punto de corte de cada par de rectas:\begin{tasks}[after-item-skip=6pt](2)
  \task $\left\{\begin{array}{l}x+y-5=0\\ 2x-y-1=0\end{array}\right.$
  \task $\left\{\begin{array}{l}3x-2y-4=0\\ x+4y+8=0\end{array}\right.$
  \task $\left\{\begin{array}{l}2x+3y=0\\ x-y-5=0\end{array}\right.$
  \task $\left\{\begin{array}{l}4x-2y+3=0\\ 2x+2y-3=0\end{array}\right.$
\end{tasks}

\ex{1} Estudia la posición relativa de estos pares de rectas comparando sus coeficientes:\begin{tasks}[after-item-skip=6pt](2)
  \task $\left\{\begin{array}{l}2x-y+3=0\\ 2x-y+3=0\end{array}\right.$
  \task $\left\{\begin{array}{l}x+3y-2=0\\ 2x+6y+1=0\end{array}\right.$
  \task $\left\{\begin{array}{l}3x-y+1=0\\ x+2y-5=0\end{array}\right.$
  \task $\left\{\begin{array}{l}\tfrac12x+y-1=0\\ x+2y-2=0\end{array}\right.$
\end{tasks}

\ex{1} Estudia analíticamente la posición relativa de los siguientes pares de rectas. Si se cortan, calcula el punto de corte:\begin{tasks}[after-item-skip=6pt](2)
  \task $r:2x+3y-5=0$\\ $s:2x-3y-11=0$
  \task $r:2x-y-3=0$\\ $s:-2x+y-1=0$
  \task $r$: recta que pasa por $A(2,3)$ y $B(4,7)$\\ $s$: recta que pasa por $C(2,7)$ y $D(7,8)$
  \task $r$: recta que pasa por $A(2,4)$ y $B(4,7)$\\ $s:y=\frac32x-2$
\end{tasks}

\ex{2} Estudia la posición relativa de las rectas y, si se cortan, halla el punto de corte:\begin{tasks}[after-item-skip=8pt](2)
  \task $r:(x,y)=(1,2)+\lambda(2,-1)$\\ $s:\dfrac{x-3}{-4}=\dfrac{y-1}{2}$
  \task $r:\left\{\begin{array}{l}x=2+\lambda\\ y=1-3\lambda\end{array}\right.$\\ $s:(x,y)=(0,4)+\mu(2,-6)$
  \task $r:\left\{\begin{array}{l}x=-1+2\lambda\\ y=3+\lambda\end{array}\right.$\\ $s:x-y+1=0$
  \task $r:2x+y-4=0$\\ $s:(x,y)=(1,1)+\lambda(1,-1)$
\end{tasks}

\ex{2} Determina el punto de corte de las rectas $r:\dfrac{x-1}{2}=\dfrac{y+2}{3}$ y $s:y-1=-(x-4)$, pasando antes las dos a forma general.

\ex{2} Parámetros:\begin{enumerate}[label=\alph*),leftmargin=*,itemsep=2pt,topsep=2pt]\item ¿Cuánto tiene que valer $a$ para que las rectas $r:y=ax-3$ y $s:\dfrac{x-1}{3}=\dfrac{y+2}{-6}$ sean paralelas?\item ¿Cuánto tiene que valer $b$ para que $3x+by-2=0$ y $6x-4y+1=0$ sean paralelas? ¿Pueden ser coincidentes para algún valor de $b$?\end{enumerate}

\ex{3} Problemas con parámetros:\begin{enumerate}[label=\alph*),leftmargin=*,itemsep=2pt,topsep=2pt]\item Calcula $a$ y $b$ para que las rectas $ax+3y+b=0$ y $2x-y+4=0$ sean coincidentes.\item Halla $m$ para que las tres rectas $x+y-3=0$, $2x-y=0$ y $mx+2y-5=0$ pasen por un mismo punto.\end{enumerate}


% =====================================================
\seccio{8}{Rectas paralelas y perpendiculares. Mediatriz}

\exf{1} Escribe la ecuación de la recta $s$, paralela a $y=\frac12x$, que pasa por el punto $A(4,4)$. Calcula el punto de corte de la recta $s$ con la recta $r$ del ejercicio \ref{ex:r7}.

\ex{1} Calcula la ecuación general de la recta perpendicular a $3x-y=1$ que pasa por el punto $A(0,1)$.

\ex{1} Calcula la ecuación de la recta paralela a la recta $2x-3y+4=0$ que pasa por el punto $A(-1,2)$.

\ex{2} Dada la recta $r:3x+y-2=0$, calcula la recta $s$ paralela a $r$ y la recta $t$ perpendicular a $r$ que pasan ambas por el punto $P(2,-1)$. Calcula el punto de corte de las rectas $r$ y $t$.

\ex{1} Calcula la ecuación de la recta mediatriz del segmento formado por los puntos $A(4,0)$ y $B(-2,6)$.

\ex{1} Calcula las ecuaciones de una recta paralela y de una perpendicular a cada una de estas rectas que pasen por el punto $P(-2,3)$:\begin{tasks}(4)
  \task $2x+5y-1=0$
  \task $x-4y+3=0$
  \task $-3x+y+2=0$
  \task $6x+2y-7=0$
\end{tasks}

\ex{1} Rectas en forma explícita:\begin{enumerate}[label=\alph*),leftmargin=*,itemsep=2pt,topsep=2pt]\item Halla la ecuación de la recta paralela a $y=-3x+1$ que pasa por el punto $B(2,-4)$.\item Determina la ecuación de la recta que pasa por el punto $A(4,1)$ y es perpendicular a la recta $y=4x-5$.\end{enumerate}

\ex{2} Escribe las ecuaciones pedidas:\begin{enumerate}[label=\alph*),leftmargin=*,itemsep=2pt,topsep=2pt]\item Ecuaciones paramétricas de la recta paralela a $r:\left\{\begin{array}{l}x=3+2t\\ y=-1-5t\end{array}\right.$ que pasa por $C(1,4)$.\item Ecuación continua de la recta perpendicular a $\dfrac{x+1}{4}=\dfrac{y-2}{-1}$ que pasa por $B(0,-3)$.\end{enumerate}

\ex{2} Rectas que pasan por un punto de corte:\begin{enumerate}[label=\alph*),leftmargin=*,itemsep=2pt,topsep=2pt]\item Calcula la ecuación de la recta paralela a $\dfrac{x-2}{3}=\dfrac{y}{-2}$ que pasa por la intersección de las rectas $y=2x-1$ y $x+y-5=0$.\item Escribe la ecuación de la recta perpendicular a $\left\{\begin{array}{l}x=1+2t\\ y=-t\end{array}\right.$ que pasa por el punto de intersección de las rectas $x-y+1=0$ y $2x+y-7=0$.\end{enumerate}

\ex{2} Considera la recta $r:y=-\frac23x+2$.\begin{enumerate}[label=\alph*),leftmargin=*,itemsep=2pt,topsep=2pt]\item Calcula la ecuación explícita de la paralela a $r$ que pasa por $P(3,-1)$.\item Determina la ecuación punto-pendiente de la perpendicular a $r$ que pasa por el origen.\item Escribe en forma continua la paralela a $r$ que pasa por $Q(0,5)$.\item Calcula la paralela y la perpendicular a la recta $y=-4$ que pasan por $(3,2)$.\end{enumerate}

\ex{2} Calcula la mediatriz del segmento de extremos $A(-3,1)$ y $B(5,-3)$. Comprueba que el punto $P(3,3)$ está en la mediatriz y que está a la misma distancia de $A$ que de $B$.

\ex{3} En el triángulo de vértices $A(1,1)$, $B(7,3)$ y $C(3,7)$, calcula la ecuación de la altura que sale de $C$ (la perpendicular a $AB$ que pasa por $C$), el pie $H$ de esta altura y el área del triángulo.

\ex{3} Los puntos $A(-1,0)$ y $B(3,2)$ son dos vértices consecutivos de un rectángulo $ABCD$. El vértice $C$ está sobre la perpendicular a $AB$ que pasa por $B$ y tiene abscisa $x=1$. Calcula $C$ y $D$, y comprueba que, de hecho, $ABCD$ es un cuadrado.

\vfill
\begin{center}\color{muted}\small\faCheckCircle[regular]\ \ ¡Ánimo! Comprueba tus resultados con las herramientas interactivas de los apuntes de la unidad.\end{center}

\end{document}
